\section{The corrected unmarked alphabet on the actual block}
\label{sec:actual-corrected-body-compiler}

The marked expansion needs one unmarked alphabet chosen before the
number and locations of the insertions. We construct that alphabet
from the actual block exponent. Its reconstruction identity concerns
the original fine density. Its smallness and normalization estimates
also use the representation with the additive remainder removed.
Keeping these two uses on the same alphabet is essential.

\subsection{The fixed data and the remaining analytic inputs}

Fix one step with fine and coarse tori, its actual path pair, clips
and conditional Gaussian construction. Write $P$ for the fine
representation, $\beta>0$ for its coupling, $\theta$ for its flag
threshold, and $J$ for the local Haar chart ratio. Fix a reference
forest $F_0$ and tag $k_0$. Put $y=r(x)$ for the recording
parameter, $B=\lceil(\log y)^2\rceil$, and $z=\beta_F(x)$
for the adjacent fine parameter. All the objects below use this
one choice of data. In particular, the fine path and clip in the
Gaussian construction are exactly those of $P$.

The following hypotheses specify the analytic interface of this
section. The actual Gaussian data have the previously established
positive floors and local decompositions; their two determinant
series and inverse correction have the purchased support prices.
The Gaussian coordinate and linear-residual bounds on the chart
mask are respectively $y^2$ and $y^{-4}$. The coarse budget
radius $B_c$ and fine radius $B_f$ satisfy
\[
 25B+30\le B_c\le B^2,\qquad B_f\le2B_c.
\]
The inherited activity estimate has its positive support-exponent
reserve, the actual path-counting bound, and
\[
 \Gamma\le\beta_F(y)^{1/8},\qquad
 \delta\le y^{-1/2},\qquad
 \beta\le4y,\qquad \theta\le2\,\mathrm{fiery}(y).
\]
The fine coupling and threshold also obey
$z/2\le\beta\le2z$ and $\theta=z^{-1/2}\mathrm{hot}(z)$.
The coarse budget threshold is exactly $\operatorname{temp}(y)$,
the fine threshold satisfies $\theta\ge\mathrm{fiery}(y)/2$,
and the profile width $W_{\rm prof}$ is positive.
The chart radius is a fixed nonnegative multiple of
$\mathrm{fiery}(y)$, and its cutoff and logarithmic radius
conditions hold. These are actual data conditions, uniform in
period, forest and tag; they cannot be supplied by choosing an
unrelated representation with the same nominal bounds.

Two further inputs must be distinguished. The owner-energy
estimate bounds the harvested ceiling charge for the actual seeded
packet whenever its chart doors are nonzero. The warning estimate
is uniform over all compatible field/store pairs: it bounds the
actual Gaussian palette and makes the warning factor vanish under
the stated local field and Gaussian bounds. A comparison on posed
reader patches alone is insufficient for this second assertion.
Their exact applications to the harvested packets and warning
factors remain prerequisites of the present construction.

\subsection{A local decomposition of the real exponent}

An admissible configuration of the local calculation consists of a
store, forest, tag, coarse field, frozen links and Gaussian
coordinates. Denote it by $\rho$. Compatibility means that the
stored coarse flags agree with those of the coarse field. Let
$\mathcal A(\rho)$ be the event that the product of the chart,
profile, likelihood and wire masks is nonzero. Its factors are
the actual factors in the block integral, including the finite tag
count in the site and bond factors. Precisely,
\[
 \prod_j\mathrm{mask}_j
 =\mathrm{profile}\cdot\#\mathrm{Choices}\cdot
      \mathrm{likelihood}\cdot\mathrm{doors}.
\]
This tag multiplicity is present in the pointwise sliced identity;
it cancels the denominator of the normalized tag average when the
block integral is reconstructed.

Let $\mathcal O(\rho)$ be the finite collection of owners of
the seeded components. The real exponent remaining after extraction
of the scalar Gaussian normalization decomposes as
\begin{equation}\label{eq:actual-body-exponent-split}
 E(\rho)=\sum_{i\in\mathcal O(\rho)}
                \bigl(c_i(\rho)+r_i(\rho)\bigr)
       +T(\rho).
\end{equation}
Here $c_i$ is minus $\beta$ times the actual ceiling charge
assigned to owner $i$, and $r_i$ is its allocated local
determinant and frozen-coordinate normalization. The remaining
term is the sum of eight specified contributions: exterior
ceiling charges, the linear residual, the quadratic inverse
correction, exterior local determinants, the relative logarithmic
determinant series, its inverse-correction logarithmic series,
the local Haar logarithm, and the inherited fine/coarse remainder.
This identity is obtained by partitioning the ceiling charges into
owner ceilings and their complement, partitioning the local
determinants in the same way, and retaining both convergent
determinant series. The linear and Haar terms are sums over the
active coordinates; the inverse correction is its absolutely
convergent path expansion. These operations are justified by the
prices below and the determinant identities of
Section~\ref{sec:gaussian-reassembly}.

For clarity, a local price at radius $B_c$ and exponent $e\ge0$
means nonnegative numbers $b_j$ and nonempty supports $X_j$ with
\begin{equation}\label{eq:actual-body-local-price}
 \sup_v\sum_{j:v\in X_j}
       e^{e\,\operatorname{cov}_{B_c}(X_j)}b_j\le t,
 \qquad |L_j(\rho)|\le b_j\quad\hbox{on }\mathcal A.
\end{equation}
The functions $L_j$ depend measurably only on the indicated local
data. The cover number is the one used for the inherited activity
prices. Set
\[
 q=y^{-1/8},\qquad \nu=y^{-5},\qquad a_y=y^{-1/24}.
\]
The actual decomposition has the following prices:
\[
\begin{array}{l|c}
 \text{contribution to }T&\text{local price}\\ \hline
 \text{exterior ceiling charges}&q\\
 \text{linear residual}&q\\
 \text{quadratic inverse correction}&y^4\nu\\
 \text{exterior local determinants}&q\\
 \text{relative logarithmic determinant series}&\nu\\
 \text{inverse-correction logarithmic series}&\nu\\
 \text{local Haar logarithm}&q\\
 \text{inherited fine/coarse remainder}&a_y
\end{array}
\]

We spell out the conversion of the earlier estimates into this
table. Each exterior ceiling term is at most
$4\,\mathrm{Aff.weight}\,y^{-1/4}$ on $\mathcal A$.
Its support lies in a coarse halo of radius $420B$ at its
assigned address. The bounded address density and halo volume
cost only a fixed polynomial in $B$; the cover number at radius
$B$ is bounded by a fixed constant. A polynomial in $B$ times
$y^{-1/4}$ is eventually at most $y^{-1/8}$. This proves the
first line. For the linear term, the product of a residual
coordinate bounded by $y^{-4}$ and a Gaussian coordinate
bounded by $y^2$ is at most $y^{-2}$. Sum over the fixed color
dimension and use the same halo calculation to obtain the second
line. The quadratic term has two such Gaussian factors, so the
inverse-correction path price is multiplied by at most $y^4$.
The factor $1/2$ in the quadratic form only improves the bound.

Outside the owner neighbourhood of radius $200B$ there is no
frozen-coordinate charge. The actual local determinant estimate
there is at most $y^{-1/4}$, and its radius-$420B$ support
gives the fourth line by the first calculation. The two remaining
determinant series have their actual word supports and prices
$\exp(-dB)$ with fixed positive decay rates $d$. Since
$B\ge(\log y)^2$, each is eventually at most $y^{-5}$:
it suffices that $d\log y\ge5$. This gives the next two
lines without discarding either series.

The extended Haar logarithm satisfies
$|\log J_{\rm ext}(u)|\le C\min(\|u\|,4\,\mathrm{arc})^2$.
The chart radius assumption and
$\mathrm{fiery}(y)^2\le y^{-1/4}$ therefore give another
$Cy^{-1/4}$ bound on the same local halo, proving the seventh
line. This estimate remains valid for the extended chart, so no
unproved restriction on $u$ is inserted.

Finally apply Section~\ref{sec:inherited-activity-prices} with
the actual positive reserve and $B_f\le2B_c$. The inherited
label term costs at most a fixed multiple of $B^8\delta$,
which is at most $y^{-1/8}$ eventually. The two fine/coarse
tail terms each cost $k l\,\beta_F(y)^{-1}$, and the
quadratic, cubic and quartic path terms each cost
$k\beta_F(y)^{-1/12}$. Here $k$ is the number of signed
orientations and $l$ the number of plaquette axes. The path
coefficients have the required smallness because
\[
 |\mathrm{Aff.weight}|\theta^2\Gamma,\quad
 |\beta\,\mathrm{Aff.weight}|\theta^3 C,\quad
 |\beta\,\mathrm{Aff.weight}|\theta^4 C
 \le \beta_F(y)^{-1/8}
\]
eventually, with the fixed path bound $C$ and the hypotheses
above. This is the specialization to the quadratic, cubic and quartic terms
of the inherited estimate. Since $\beta_F(y)\ge y$, its total
price is at most
$(1+2kl+3k)y^{-1/12}\le y^{-1/24}$ eventually.
Increasing the support radius from $B$ to $B_c$ decreases its
cover number and preserves every preceding price.

Thus $T$ has total price
\begin{equation}\label{eq:actual-body-price-total}
 t_y=q+q+y^4\nu+(q+(\nu+\nu))+q+a_y.
\end{equation}
For $y\ge1$, $\nu\le y^4\nu\le q\le a_y$, whence
$2t_y\le16a_y$. Eventually $16y^{-1/24}\le y^{-1/48}$.
In particular $t_y\le1$ and $2t_y\le y^{-1/48}$.

\subsection{Bounded extension and averaging}

Let $s_i$ be the number of seed shares belonging to owner $i$.
Choose $\alpha>0$ with $2\alpha\le a/2$, where $a>0$
is the owner-energy constant. That estimate and the local
normalization bound give, on $\mathcal A$ for occurring owners,
\[
 c_i\le-\beta a\theta^2s_i,\qquad
 |r_i|\le\alpha\,\mathrm{hot}(z)^2s_i.
\]
Indeed the normalization is bounded by
$C(1+B)^4\log(y)\,s_i$ with a fixed constant $C$.
The scale payment makes this at most
$\alpha\,\mathrm{hot}(z)^2s_i$ eventually.
Since $z\theta^2=\mathrm{hot}(z)^2$ and $\beta\ge z/2$,
\[
 2\alpha\,\mathrm{hot}(z)^2\le\beta a\theta^2.
\]
Consequently $c_i+r_i\le-A_i$, where
$A_i=\alpha\,\mathrm{hot}(z)^2s_i$.
The finite-volume lower estimate for the same ceiling charges
also gives $c_i\ge u$ with finite $u$, uniformly over the
compatible data on $\mathcal A$. It need not be uniform in the
volume. Set
\[
 D_i=\min(-A_i,u-A_i),\qquad
 b_i(\rho)=\max\{D_i,\min\{-A_i,c_i(\rho)+r_i(\rho)\}\}.
\]
Then $D_i\le b_i\le-A_i$ everywhere on compatible data,
and $b_i=c_i+r_i$ for occurring owners on $\mathcal A$.
The clamping operation is continuous, preserves locality and
makes $b_i$ bounded. This extension permits averaging before
the nonzero-mask event is imposed.

Similarly replace every local summand $L_j$ of $T$ by its clamp
to $[-b_j,b_j]$, using the prices in
\eqref{eq:actual-body-local-price}. This changes no value on
$\mathcal A$. Coalesce terms with the same nonempty support.
The sums converge absolutely: summing the site prices over the
finite torus bounds $\sum_j b_j$, because each support contains
a site and $e\ge0$. The grouped local terms $h_X$ retain the
same site price and sum to $T$ on $\mathcal A$.

Average both the owner functions and the grouped terms over the
finite signed lattice motions and the product of local Haar
gauge measures. Transport supports, owners, stores, forests and
tags together. The actual exponent and the nonzero-mask event are
invariant under this common action. Thus averaging
\eqref{eq:actual-body-exponent-split} gives local invariant
functions $\bar b_i,\bar h_X$ satisfying
\begin{equation}\label{eq:actual-body-averaged}
 E=\sum_{i\in\mathcal O}\bar b_i+\sum_X\bar h_X
 \quad\hbox{on }\mathcal A,\qquad
 \bar b_i\le-\alpha\,\mathrm{hot}(z)^2s_i.
\end{equation}
Boundedness makes all these averages integrable. The action
preserves seed shares and the cover number, so the averaged
majorants $Y_X$ satisfy the same site price $t_y$.
For any site of $X$, that price gives $Y_X\le t_y\le1$.
Therefore
\[
 |e^{\bar h_X}-1|\le2|\bar h_X|\le2Y_X.
\]
The auxiliary letters $e^{\bar h_X}-1$ have total local price
at most $2t_y\le y^{-1/48}$. Together with the negative owner
functions, these are the corrected unmarked alphabet. This is a
value estimate; it does not assert arbitrary differentiated
estimates merely from bounded extension.

\subsection{The scalar normalization and the genuine density}

Write $K_0$ for the bare kinetic matrix at the fixed reference
forest $F_0$ and tag $k_0$, $N_f$ for the number of fine sites,
and $n_0$ for the number of free reference bonds. The Lie algebra
dimension is already included in the factor $\mathrm{spare}$.
The
scalar extracted from every actual sector is
\begin{equation}\label{eq:actual-body-loading}
 \Lambda(V)=\mathrm{bank}(V)\,\mathrm{pref}(\beta)
       e^{P.v N_f}\,\mathrm{spare}(J,\beta)^{n_0}
       (\det K_0)^{-1/2}>0.
\end{equation}
All factors are positive, including the determinant by the bare
kinetic floor. The forest determinant and free-bond count
identities of Section~\ref{sec:gaussian-reassembly} identify
this same scalar for every actual forest and tag. No new scalar
is selected after a marked family is specified.

Here is the exponent check. Let $Q$ be the actual covariance,
$W$ its Gram precision, $m$ its localized mean, $Z$ its Gaussian
coordinate and $g=\beta^{-1/2}$. Let $U$ be the fine field
constructed from $m+gZ$ and the frozen links. Denote the
quadratic guide by $G_0$, the fine cut action by $R_f$, the
coarse quadratic action by $e_0$, and the prediction mismatch
cost by $M_0$. The charge decomposition and quadratic completion
give respectively
\begin{align*}
 \mathrm{Aff.weight}\sum\mathrm{levy}
   &=M_0+(R_f-e_0)-G_0(m+gZ)+G_0(m),\\
 -\beta G_0(m)+\eta\cdot Z
       +\tfrac12 Z^\top(Q^{-1}-W)Z
   &=\tfrac12\mathrm{force}_Q(Z)-\beta G_0(m+gZ).
\end{align*}
The inherited remainder satisfies
$P.vN_f+\mathrm{apart}(U)=P.\mathrm{exponent}(U)+\beta R_f$.
The bank identity is
$e^{-\beta e_0}=\mathrm{bank}(V)e^{\mathrm{tail}_c(V)}$.
Substitution into the definition of $E$ cancels $R_f$ and both
guide terms and yields
\[
 E=E_{\det}+\mathrm{tail}_c+\tfrac12\mathrm{force}_Q(Z)
       +\log J_{\rm ext}+\beta e_0-\beta M_0
       +P.\mathrm{exponent}(U)-P.vN_f.
\]
Here $E_{\det}$ includes the frozen-coordinate logarithm and
both determinant logarithms. Precisely,
\[
 J.\mathrm{scale}^{\#\mathrm{Active}}\,
       \mathrm{chartCost}(Q,\beta^{-1/2})
 =\mathrm{spare}(J,\beta)^{n_0}(\det K_0)^{-1/2}e^{E_{\det}},
\]
as in \eqref{eq:reassembly-overhead}--\eqref{eq:reassembly-logs};
$\#\mathrm{Active}$ also counts bonds.
Also the actual prediction noise equals
$\mathrm{pref}(\beta)e^{-\beta M_0}$, and the product of
extended Haar factors is $e^{\log J_{\rm ext}}$. Multiplying
these identities proves the pointwise identity
\begin{equation}\label{eq:actual-body-density}
 \mathrm{freight}(\rho)=\Lambda(V)e^{E(\rho)},
\end{equation}
where freight is the original fine exponential times the
prediction noise, chart Jacobian, Gaussian density conversion
and extended Haar factors. Thus the exponential here is that
of the genuine $P$.

\subsection{Flags, connected components and marked reuse}

Expand each auxiliary exponential as $1+(e^{\bar h_X}-1)$.
Outside the owner regions the warning letter is the corresponding
mask minus one; hence its factor $1+\mathrm{warning}$ is
exactly that mask. Within an owner region, the owner letter
already includes the product of these masks. Owner regions are
disjoint. Multiplying the factors therefore recovers
$e^{\sum\bar b_i+\sum\bar h_X}$ times the product of all
actual masks.

There is also no multiplicity from decorating owners with fine
flags. Each owner decoration is required to equal the fine flags
in its assigned region, so exactly one decoration survives for
each owner. On $\mathcal A$, the warning/door implication places
every actual fine flag in an owner region. The union of these
unique decorations is therefore exactly the actual fine flag set.
Grouping overlapping letters into connected components gives a
bijection between compatible component collections and the
underlying letter choices, as in
Section~\ref{sec:allarity-planted-compiler}. It preserves both
flag unions and the stored seed data. Consequently the sum with
fine flag inventory $u$ is
\[
 \mathbf1_{\{\mathrm{flags}_\theta(U)=u\}}
       e^{E(\rho)}\prod_j\mathrm{mask}_j(\rho).
\]
Outside $\mathcal A$ both sides vanish because the full mask
product is zero; no division by a mask is used. Combining this
identity with \eqref{eq:actual-body-density} proves the genuine
body identity \eqref{eq:planted-body} for this corrected alphabet
and the scalar \eqref{eq:actual-body-loading}.

For the same alphabet, apply the real additive correction and
masked logarithm construction of
Section~\ref{sec:restarted-rank-history}. The required warning
and product prices, support reserve and signed/gauge transport
must be those of the alphabet just constructed. With these
hypotheses the all-hole normalization and forced-exclusion
identities of Section~\ref{sec:allarity-planted-compiler} apply
to the white representation. They retain the same scalar and
coarse output chosen by the masked logarithm construction.
Only then introduce an arbitrary finite family of marked sites
and its parent map. The marked decorations vary with that family;
the body alphabet, genuine body identity and scalar do not.
Decreasing the already purchased exponents and multipliers gives
both uses in Section~\ref{sec:actual-mixed-marked-compiler}.
The separate harvested-energy and all-compatible-warning
prerequisites at the start remain necessary for this application.
